\documentclass[12pt]{article}
\usepackage[utf8]{inputenc}
\usepackage[T1]{fontenc}
\usepackage{hyperref}
\title{Yareni Zapotec Language Resources for K-12 Curriculum Development}
\author{Compiled by Claude Sonnet 5.5 (claude-sonnet-5-5) for Andr\'as Kornai}
\date{}
\begin{document}
\maketitle

\section{Introduction}
Yareni Zapotec (ISO 639-3 zae, Glottocode yare1249) is spoken in Santa Ana Yareni, Oaxaca, Mexico, by about 2,900 people. It belongs to the Ixtl\'an Zapotec cluster. Little material exists. No dedicated dictionary or school materials were found.

\section{Primary Educational Resources}
No primers, textbooks or school curricula were found in the search results.

\section{Linguistic and Reference Materials}
\begin{itemize}
\item Hungarian HLT database entry, which lists several online descriptions and resources: \url{https://hlt.bme.hu/en/dld/language/2990}
\item Wikipedia article on Ixtl\'an Zapotec, the related cluster. It cites older materials on related Sierra Ju\'arez varieties: \url{https://en.wikipedia.org/wiki/Ixtl%C3%A1n_Zapotec}
\item Wikitongues language page: \url{https://wikitongues.org/languages/zae/}
\item Number-counting page for Santa Ana Yareni Zapotec: \url{https://www.languagesandnumbers.com/how-to-count-in-santa-ana-yareni-zapotec/en/zae/}
\end{itemize}

\section{Religious and Community Resources}
\begin{itemize}
\item Scripture Earth page for the Yareni Zapotec New Testament: \url{https://www.scriptureearth.org/00eng.php?iso=zae}
\item Digital Bible Society edition of the Yareni Zapotec Bible: \url{https://dbs.org/bibles/ZAEWBT}
\end{itemize}

\section{Digital and Audio Resources}
\begin{itemize}
\item Global Recordings Network audio and slideshows in Yareni Zapotec: \url{https://www.globalrecordings.net/en/language/zae}
\end{itemize}

\section{Curriculum Development Context}
The language is thought to be used by the whole community, but it is not known to be taught in schools. The New Testament text and the audio could serve as sources of written and spoken material. Sources disagree on vitality.
\end{document}